\subsection{Nilpotent Rees quotients with arbitrary row and column actions}

\begin{proposition}\label{E:band-free-group-bar}
Let $M$ be a finite monoid whose minimal ideal is a rectangular band
$K=I\times J$.  Write $Q=M/K$ for the pointed Rees quotient, with the
entire ideal identified to its zero.  For a finite nonempty set $I$
with chosen element $0$, let $F_I$ be the free group on symbols $x_i$
for $i\ne0$, and set $x_0=1$.  A set map $\sigma:I\to I$ induces the
group homomorphism
\[
 \Phi_\sigma(x_i)=x_{\sigma(0)}^{-1}x_{\sigma(i)}.       \tag{E.26}
\]
The row actions of $M$ on $I$ give a pointed left $Q$-action on $BF_I$,
and the column actions give a pointed right $Q$-action on $BF_J$.
There is a stable equivalence
\[
 \widetilde{\Sigma^\infty BM}\simeq
 \Sigma^\infty B(BF_J,Q,BF_I).                         \tag{E.27}
\]
Here the bar is the \emph{pointed, smash-product} two-sided bar: its
degree-$n$ space is $BF_J\wedge Q^{\wedge n}\wedge BF_I$.
\end{proposition}

\begin{proof}
Use ordinary monoid multiplication, so $m(i,j)=(\sigma_m(i),j)$
and $(i,j)m=(i,\tau_m(j))$.  Thus
$\sigma_{mn}=\sigma_m\sigma_n$ and
$\tau_{mn}=\tau_n\tau_m$.  Substitution in (E.26) gives the literal
identity $\Phi_\sigma\Phi_\rho=\Phi_{\sigma\rho}$:
\[
 \Phi_\sigma(x_{\rho(0)}^{-1}x_{\rho(i)})
 =x_{\sigma\rho(0)}^{-1}x_{\sigma\rho(i)}.
\]
A constant set map induces the trivial group homomorphism.  Every
element of the rectangular band acts constantly on each of $I,J$,
so these pointed actions factor strictly through $Q$.

We give the coherence in the stable identification explicitly.
Put $R=\Sigma^\infty M_+$, $A=\Sigma^\infty Q$, and choose
$e=(0,0)\in K$.  The ideal spectrum is literally
\[
 \Sigma^\infty K_+\cong Re\wedge_{\mathbb S}eR
\]
as an $R$-bimodule, by the bijection
$((i,0),(0,j))\mapsto(i,j)$.  Since $Re,eR$ are idempotent retracts
of regular modules and $eRe=\mathbb S$, the ideal--quotient cofiber
sequence gives
$A\wedge_R^{\mathbb L}A\simeq A$.  Restriction of $A$-modules to
$R$-modules is therefore fully faithful.

Let $V=\operatorname{fib}(Re\to\mathbb S)$ and
$W=\operatorname{fib}(eR\to\mathbb S)$ for the trivial augmentation
modules.  Applying $e$ gives the fiber of the identity of $\mathbb S$,
so $eV=We\simeq0$ as spectra.  Tensoring the ideal--quotient sequence
then supplies coherent equivalences
$V\simeq A\wedge_R^{\mathbb L}V$ and
$W\simeq W\wedge_R^{\mathbb L}A$.  The two augmentation fiber
sequences, exactly as in (E.25), give
\[
 \widetilde{\Sigma^\infty BM}
 \simeq\Sigma^2(W\wedge_A^{\mathbb L}V).               \tag{E.28}
\]

Let $\Gamma_I$ be the graph with two vertices $v_0,v_1$ and one
oriented edge $a_i:v_0\to v_1$ for every $i\in I$, based at $v_1$.
The cofiber of the fold $\Sigma^\infty I_+\to\mathbb S$ is naturally
$\Sigma^\infty\Gamma_I$.  Consequently
$\Sigma V\simeq\Sigma^\infty\Gamma_I$, naturally for all set maps
of $I$.  The same statement holds for columns.

For precision one can implement the next natural equivalence by a
literal equivariant zigzag.  Let $\mathcal G_I$ be the free groupoid
on the directed graph $\Gamma_I$.  There are natural based maps
\[
 \Gamma_I\longrightarrow B\mathcal G_I
       \longleftarrow B\operatorname{Aut}_{\mathcal G_I}(v_1).
                                                               \tag{E.29}
\]
The first map sends an edge to its generating arrow.  The second is
the full subgroupoid inclusion at $v_1$.  Both are homotopy
equivalences: the graph is aspherical and has this fundamental
groupoid, and $\mathcal G_I$ is connected.  Both maps commute with
every set map $I\to I$, because such maps fix the two vertices.
Identifying $x_i=a_0^{-1}a_i$ identifies the last group with $F_I$
and its induced maps with (E.26).  Thus (E.29), not merely a
degree-zero module calculation, gives an $R$-module equivalence
\[
 \Sigma V\simeq\Sigma^\infty BF_I.
\]
The right side has the strict $A$-module structure already described;
full faithfulness of restriction upgrades this to an equivalence of
$A$-modules.  Similarly $\Sigma W\simeq\Sigma^\infty BF_J$ as right
$A$-modules.  Substitute in (E.28).  Suspension spectra preserve
smash products and geometric realizations, so the derived tensor
product of these two suspension spectra is the suspension spectrum
of the pointed two-sided bar.  This proves (E.27).
\end{proof}

\begin{lemma}\label{E:free-subset-intersections}
For $S\subseteq I$, put
$H_S=\langle x_a^{-1}x_b:a,b\in S\rangle\leq F_I$,
with $H_\varnothing=1$.  Then
\[
 \operatorname{im}\Phi_\sigma=H_{\sigma(I)},\qquad
 H_S\cap H_T=H_{S\cap T}.                              \tag{E.30}
\]
Every finite union of the nerve subcomplexes $BH_S\subset BF_I$
has finite CW homotopy type.  The same holds for finite unions of
products or smash products of such subcomplexes, and for cofibers of
inclusions between these unions.
\end{lemma}

\begin{proof}
The first identity follows directly from (E.26): its generators are
$x_{\sigma(0)}^{-1}x_b$ for $b\in\sigma(I)$, and their pairwise
quotients generate all $x_a^{-1}x_b$.  Regard $H_S$ as the image of
the fundamental group, based at $v_1$, of the subgraph with the
edges indexed by $S$.  A reduced edge loop in a graph is unique.
If the same element has reduced representatives using only $S$ and
only $T$, its reduced representative uses only $S\cap T$.
This proves the intersection identity, including the empty case.
For nonempty $S$, the group $H_S$ is free of rank $|S|-1$.

Nerve intersections are literal subgroup intersections, hence any
finite intersection of the displayed subcomplexes is the classifying
space of a finitely generated free group, and has finite graph
homotopy type.  To pass from intersections to unions, induct on the
number of subcomplexes.  The union of two CW subcomplexes is their
homotopy pushout over their intersection; the latter is itself a
union of fewer such intersections.  Homotopy pushouts of spaces of
finite CW homotopy type have finite CW homotopy type, by cellular
approximation of the two maps between finite replacements.
This proves the assertion for unions.  Finite products and smash
products have the same property, and their intersections have the
corresponding product form.  The same induction applies to them.
Finally, the indicated inclusions are CW cofibrations, so their
quotients are homotopy cofibers of maps between spaces of finite
CW homotopy type.
\end{proof}

\begin{theorem}\label{E:nilpotent-Rees-quotient}
Let $M$ be a finite monoid with rectangular-band minimal ideal $K$.
Suppose its pointed Rees quotient is $Q=N^1$, where $N$ is a finite
nilpotent semigroup with zero and the identity is freshly adjoined.
No fixed-row or fixed-column hypothesis is imposed on the outside
actions.  Then, for every prime $p$, $H^*(BM;\mathbb F_p)$ is locally
finite over the Steenrod algebra.  In particular $BM\not\simeq
\Sigma\CP^\infty$.
\end{theorem}



\begin{proof}
Write $X=BF_I$, $Y=BF_J$ with their pointed actions from
Proposition~\ref{E:band-free-group-bar}, and put
$Z=B(Y,Q,X)$.  Choose $\ell\geq2$ with $N^\ell=\{0\}$ and set
\[
 w(a)=\max\{r:a\in N^r\},\qquad a\ne0.
\]
Thus $1\leq w(a)\leq h:=\ell-1$, and nonzero products satisfy
$w(ab)\geq w(a)+w(b)$.  Define descending based subcomplexes
\[
 X^0=X,\qquad
 X^s=\bigcup_{\substack{a\ne0\\w(a)\geq s}}
          B\operatorname{im}\Phi_a\quad(s\geq1),       \tag{E.31}
\]
where an empty union means the basepoint.  Define $Y^s$ in the same
way for the right action.  In particular $X^{h+1}=Y^{h+1}=*$.
For a nonzero nonunit $b$, the action maps satisfy
\[
 bX^s\subseteq X^{s+w(b)},\qquad
 Y^s b\subseteq Y^{s+w(b)}.                             \tag{E.32}
\]
For $s=0$ this is the definition of the image of $b$.  For $s>0$ it
follows by composing $b$ with the actor defining an image in (E.31)
and using superadditivity of $w$; a zero product gives the basepoint.
Every subcomplex in (E.31), and every finite union or intersection
needed below, has finite CW homotopy type by
Lemma~\ref{E:free-subset-intersections}.

A nondegenerate middle word in bar degree $n$ is
$\mathbf a=(a_1,\ldots,a_n)$ with all $a_i\ne0,1$; write
$w(\mathbf a)=\sum_iw(a_i)$.  In the summand $Y\wedge X$ indexed by
this word, define the level-$L$ subspace to be
\[
 \bigcup_{s+t+w(\mathbf a)\geq L}Y^s\wedge X^t,
 \qquad 0\leq s,t\leq h.                              \tag{E.33}
\]
Give the identity letter weight zero, so this definition extends to
degenerate words.  The basepoint and all these subspaces define a
based simplicial subspace $F_LZ_\bullet$ of the pointed bar.
Indeed an inner face either gives the basepoint or replaces two
letters by one of at least their total weight.  An end face applies
the deleted letter to an endpoint and, by (E.32), transfers its
weight to the endpoint filtration.  Degeneracies insert a zero-weight
identity.  Thus every simplicial operator preserves (E.33).

There are two complementary finiteness estimates.  First, (E.33)
is the basepoint whenever $(n+2)h<L$.  In a CW realization, a
nonbasepoint bar simplex in degree $n$ contributes cells only in
dimensions at least $n$.  Hence $F_LZ:=|F_LZ_\bullet|$ has no
nonbasepoint cells below degree $\lceil L/h\rceil-2$.
In particular its connectivity tends to infinity with $L$.

Second, if a nondegenerate middle word has length $n\geq L$, its
weight is at least $L$ and (E.33) contains its entire summand, using
$s=t=0$.  The quotient simplicial space therefore has no
nondegenerate middle words in degrees $n\geq L$.  In every smaller
degree there are only finitely many words.  Each remaining summand
is the cofiber of (E.33) in $Y\wedge X$.  Formula (E.33) is a finite
union of smash products of the finite unions in (E.31), and its
intersections have the same form.  By the lemma, this cofiber has
finite CW homotopy type.  The geometric realization is built in
finitely many stages from finitely many such spaces, using their
latching cofibrations and products with finite simplices.  Thus
\[
 Z/F_LZ\quad\hbox{has finite CW homotopy type}.          \tag{E.34}
\]
This is a homotopy-type assertion: the nerves $BF_I$ themselves
usually have infinitely many cells.

Fix a cohomological degree $d$.  Choosing $L$ so that
$\lceil L/h\rceil-2>d+1$ makes the quotient map induce an
isomorphism on reduced degree-$d$ cohomology.  Hence any class in
$\widetilde H^d(Z;\mathbb F_p)$ is pulled back from the cohomology
of one space (E.34), whose total cohomology is finite-dimensional.
Naturality makes every iterated Steenrod operation on that class
the pullback of the corresponding operation on this one lift.
Its entire Steenrod orbit is therefore finite-dimensional.
Proposition~\ref{E:band-free-group-bar} transfers this conclusion
to $BM$.  The explicit nonzero operation sequences on the suspended
generator of $H^2(\CP^\infty;\mathbb F_p)$, recorded in
Theorem~\ref{E:nilpotent-inflations}, exclude $\Sigma\CP^\infty$.
\end{proof}

\begin{remark}\label{E:nilpotent-Rees-scope}
This theorem extends the constant-action inflation family.  Its proof
alone does not give Steenrod local finiteness for all finite aperiodic monoids.
Nonzero idempotents in the Rees quotient prevent the positive weight
used above.  Extending the argument to further regular ideal layers
requires additional coherent control of the endpoint module spectra;
an ordinary monoid-algebra composition series does not provide it.
The strictification in
Lemma~\ref{E:regular-controlled-localization} supplies this control,
and Theorem~\ref{E:general-controlled-multibar} proves the general
result below.
\end{remark}

\begin{lemma}\label{E:relative-bar-good-prime}
Let $Q$ be a finite pointed monoid with group of units $G$, and let
$Y,X$ be pointed simplicial sets with right and left $Q$-actions.
Suppose the action of $G$ is by simplicial automorphisms.  Form the
ordinary pointed relative bar
\[
 C_n=Y\wedge_G Q\wedge_G\cdots\wedge_G Q\wedge_G X
                                                               \tag{E.35}
\]
with $n$ middle copies of $Q$.  The balancing is ordinary orbit
balancing, not derived balancing.  If $p\nmid|G|$, the canonical map
\[
 B(Y,Q,X)\longrightarrow |C_\bullet|                  \tag{E.36}
\]
is an $\mathbb F_p$-homology equivalence.  It is an actual map of
pointed spaces, and hence its cohomology isomorphism commutes with
all Steenrod operations.
\end{lemma}

\begin{proof}
Put $k=\mathbb F_p$.  Degreewise orbit balancing is compatible with
the bar faces and degeneracies, and defines (E.36).  To verify its
homology, retain the simplicial direction inside $X,Y$ and the
separate outer bar direction.  For a simplicial set $T$ with a finite
group $H$ acting simplicially, normalized chains satisfy
\[
 N\widetilde k[T/H]\cong (N\widetilde k[T])_H.          \tag{E.37}
\]
Indeed the nondegenerate simplices of $T/H$ are precisely the orbits
of nondegenerate simplices of $T$: if an orbit is degenerate, any
representative is a translate of a degenerate simplex and is itself
degenerate.  This also holds for pointed quotients.

The $n$-fold balancing in (E.35) is the orbit quotient by $G^{n+1}$,
acting by
\[
 (y,a_1,\ldots,a_n,x)\longmapsto
 (yg_0,g_0^{-1}a_1g_1,\ldots,g_{n-1}^{-1}a_ng_n,g_n^{-1}x).
\]
The natural Eilenberg--Zilber comparison for products and smash
products is equivariant.  Since $kG$ is semisimple, coinvariants are
exact and preserve this chain equivalence.  Thus the inner chains
of (E.35) are naturally modeled by
\[
 N\widetilde k[Y]\otimes_{kG} k[Q]\otimes_{kG}\cdots
       \otimes_{kG} k[Q]\otimes_{kG}N\widetilde k[X],   \tag{E.38}
\]
where $k[Q]$ is the contracted algebra, so the zero of $Q$ is zero.
The absolute bar has the same expression with tensor products over
$k$ in place of $kG$.

For completeness, the relative bar computes the same derived tensor
product as the absolute bar.  Its augmented resolution of a left
$k[Q]$-module $L$ has terms
$k[Q]\otimes_{kG}\cdots\otimes_{kG} k[Q]\otimes_{kG}L$.
The usual extra degeneracy inserting the unit contracts the augmented
complex after forgetting the left action.  Every term is projective
as a left $k[Q]$-module, since its last tensor factor is a $kG$-module
and all $kG$-modules are projective.  The same argument degree by
degree applies to bounded-below complexes of modules, as here;
the resulting double complexes are in the first quadrant.
Consequently the quotient comparison from the absolute to the
relative bar is a quasi-isomorphism.  Applied to (E.38) this proves
(E.36), including the stated naturality.
\end{proof}

\begin{theorem}\label{E:nilpotent-Rees-with-units}
Let $M$ have rectangular-band minimal ideal $K$, and suppose
$Q=M/K$ is the disjoint union of a finite group of units $G$ and a
nilpotent ideal $N$ with zero.  At every prime $p\nmid|G|$,
$H^*(BM;\mathbb F_p)$ is locally finite over the Steenrod algebra.
In particular $BM\not\simeq\Sigma\CP^\infty$.
\end{theorem}

\begin{proof}
Apply Proposition~\ref{E:band-free-group-bar} and then
Lemma~\ref{E:relative-bar-good-prime}.  It suffices to prove local
finiteness for the relative bar (E.35), with $Y=BF_J,X=BF_I$.
For $a\in N\setminus\{0\}$ use the weight
$w(a)=\max\{r:a\in N^r\}$ and give units weight zero.
Multiplication on either side by a unit preserves $N^r$, and hence
preserves $w$.  It also permutes the collection of image subgroups
in (E.31), so the endpoint filtrations $X^s,Y^s$ are $G$-invariant.
The filtration (E.33) therefore descends to the balanced quotient.
End faces still trade letter weight for endpoint filtration, and
inner faces do not decrease weight.

In the relative bar a middle unit factor is literally degenerate:
it can be absorbed into an adjacent factor by balancing and then
is the inserted identity.  The nondegenerate middle factors are
therefore exactly letters of $N\setminus\{0\}$, modulo the finite
balancing action of $G^{n+1}$.  The high-filtration connectivity
estimate in Theorem~\ref{E:nilpotent-Rees-quotient} applies unchanged,
and the quotient has no nondegenerate middle levels $n\geq L$.

Before taking the finite balancing quotient, every level quotient
has finite CW homotopy type by Lemma~\ref{E:free-subset-intersections}.
Equation (E.37) and exactness of $k[G^{n+1}]$-coinvariants show that
its orbit quotient has finite-dimensional total $k$-homology.
The quotient relative bar has only finitely many nondegenerate
middle levels.  Its skeletal homology spectral sequence consequently
shows that its total $k$-homology, and hence total $k$-cohomology, is
finite-dimensional.  No assertion of integral finite CW type for
these orbit quotients is needed.

As before, each fixed-degree cohomology class lifts to one such
high-filtration quotient.  Its Steenrod orbit is contained in the
image of the finite-dimensional cohomology of that quotient.
This proves local finiteness, and the suspended polynomial generator
excludes $\Sigma\CP^\infty$.
\end{proof}

\begin{corollary}\label{E:two-regular-classes}
Let $M$ be a finite monoid whose only regular $\mathcal J$-classes
are its minimal ideal and its group of units (allowing these to
coincide).  Let $H$ be a maximal subgroup in the minimal ideal and
$G$ the unit group.  If $p\nmid |G||H|$, then
$H^*(BM;\mathbb F_p)$ is locally finite over the Steenrod algebra.
Consequently no such monoid has classifying space
$\Sigma\CP^\infty$.
\end{corollary}

\begin{proof}
If the minimal ideal is all of $M$, the monoid is a finite group and
its positive cohomology at the indicated prime vanishes.  Otherwise
write its Rees matrix minimal ideal as $I\times H\times J$.
Collapse each minimal $H$-class, leaving all outside elements
unchanged, to obtain a monoid $\overline M$ with rectangular-band
minimal ideal.  The proof of
Lemma~\ref{lem:B-collapse-minimal-group} works over $k=\mathbb F_p$
whenever $|H|$ is invertible.  Explicitly, classwise averaging on
$kK\cong kMe\otimes_{kH}ekM$ is a bimodule idempotent.  Its kernel
$D$ is the kernel of $kM\to k\overline M$, is projective as a right
$kM$-module, and satisfies $D\otimes_{kM}k=0$.
The resolution $0\to D\to kM\to k\overline M\to0$ and
change of rings show that the actual quotient map
$BM\to B\overline M$ is a $k$-homology equivalence.  Its cohomology
isomorphism therefore preserves all Steenrod operations, not only
products.

The pointed Rees quotient $Q=\overline M/(I\times J)=M/K$ has
unit group $G$.  Its nonunits form an ideal $N$, and have no
nonzero idempotent: such an idempotent would lift to an idempotent
outside both regular classes of $M$.  A finite semigroup with zero
and no nonzero idempotent is nilpotent.  One direct proof is to take
a nonzero product longer than the number of nonzero elements.
Two of its nonzero prefixes coincide, say $a=at$ for a nonempty
intermediate word $t$.  Then $a=at^r$ for every $r$, so an
idempotent power of $t$ is nonzero, a contradiction.  Thus $N$ is
nilpotent, and Theorem~\ref{E:nilpotent-Rees-with-units} applies.
There is always a prime avoiding the two finite group orders,
whereas the suspended polynomial generator has an infinite
Steenrod orbit at every prime.
\end{proof}
